\chapter{Zusammenfassung und Ausblick}

Im Algemeinen kann man zusammenfassend sagen, dass Lokale Klassen eher weniger gebraucht werden, Anonyme Klassen aufgrund derer Spezifität häufiger und an verschiedenen Stellen, und Innere Klassen sind in bestimmten Anwendungsbereichen wie GUIs fast Standard, aber sonst auch nicht so vielsetig einsetzbar.
Insgesamt werden aber alle Konstrukte in ähnlichen Anwendungsbereichen eingesetzt, und unterscheiden sich somit maximal in der Größenordnungen des Einsatzes.
Weiterhin sind die Konzepte aus Gründen der Sichtbarkeit und Ästhetik relevant, in wenigen Fällen ist tatsächlich die Codeefizienz höher, ausgenommen sind in diesem Fall die Anonymen Klassen.
Dies kann man auch an der Tatsache erkenne, dass für die Implementierung der eingebetteten Klassen die JVM nicht verändert werden musste, da sich die Eingebetteten Klassen sich soweit nur in der Implemtierung und leicht im Verhalten von den Top-Level Klasse unterscheiden.

% Bild einbinden
%\begin{figure}[ht!]
%\begin{center}\includegraphics[width=5cm]{src/pic/logo}\end{center}
%\caption{Das SE Logo}
%\label{Logo}
%\end{figure}

\newpage
